\section{Numerical Method}

The propagation equation mixes a transverse diffraction term, a temporal dispersion term, and a nonlinear term which depends on the local field at every point in space and time. No single numerical scheme can handle all three of them at a low cost. Therefore, the code treats each mechanism with the tool best suited to it, and recombines them at every propagation step via operator splitting. This section explains the three pieces in turn, namely the radial basis used for diffraction, the time-stepping scheme that combines them, and the population equations which must be solved alongside the field.

\subsection{Quasi Discrete Hankel Transform for the Radial Coordinate}

The beam presents a cylindrical symmetry, so the transverse Laplacian $\nabla_\perp^2$ acting on a field that depends only on $r$ reduces to $\frac{1}{r}\frac{\partial}{\partial r}\big(r\frac{\partial}{\partial r}\big)$. The natural basis for this operator is not the Fourier basis but the Hankel transform of order zero, because a diffracting wave with cylindrical symmetry is a superposition of Bessel functions $J_0(\rho r)$ rather than complex exponentials. In this basis, diffraction becomes a multiplication by a phase factor, in the same way that dispersion becomes a multiplication in the temporal Fourier domain.

Since the continuous Hankel transform cannot be evaluated on a computer, the code resorts to the quasi-discrete Hankel transform. The idea consists in sampling the field not on a uniform radial grid, but at the zeros $j_1, j_2, \dots, j_N$ of the Bessel function $J_0$. On this particular grid, the forward and inverse Hankel transforms both collapse to the same linear transformation, represented by a single symmetric matrix $Y$:
\begin{equation}
Y_{mn} = \frac{2}{j_N}\,\frac{J_0(j_mj_n/j_N)}{J_1(j_m)^2}
\end{equation}
Multiplying a field sampled at the radii $r_m = j_m R/j_N$ by $Y$ gives its Hankel transform sampled at the spatial frequencies $\rho_m = j_m/R$, and multiplying by $Y$ a second time returns the original field, since $Y$ is its own inverse. This is what the function \texttt{build\_grids} constructs and what \texttt{LinearOperator\dotb half\_linear} uses, applying $Y$ before and after the diffraction phase.

One practical consequence of this grid is that the radial samples are not evenly spaced. They are densest near the axis and spread out towards the edge of the box, which fits well a beam whose interesting structure is located near $r=0$. It also implies that anything reading the radial profile afterwards must use the same non-uniform grid, which explains why the Abel matrix is built from shell integrals over these specific radii instead of a uniform trapezoidal sum.


\subsection{Strang Splitting of the Propagation Step}

The full propagation equation can be written schematically as $\partial_z u = (\mathcal{L}+\mathcal{N})u$, where $\mathcal{L}$ groups the linear operators, diffraction and dispersion, and $\mathcal{N}$ groups the nonlinear terms. Solving $\mathcal{L}$ alone and $\mathcal{N}$ alone are both cheap, since diffraction and dispersion are diagonal in the Hankel and Fourier bases, and the nonlinear term is diagonal in real space. However, solving $\mathcal{L}+\mathcal{N}$ together is not.

Strang splitting exploits this by approximating the exact evolution operator with a symmetric product:
\begin{equation}
e^{(\mathcal{L}+\mathcal{N})dz} \approx e^{\mathcal{L}\,dz/2}\,e^{\mathcal{N}\,dz}\,e^{\mathcal{L}\,dz/2}
\end{equation}
The operators do not commute, so this product differs from the exact evolution, but the error is of the third order in $dz$ per step instead of the second order for the naive product $e^{\mathcal{L}dz}e^{\mathcal{N}dz}$.

The linear half-step applies diffraction and dispersion together in a single operation rather than one after the other. This is exact and not an extra approximation, because once the field is expressed in the combined basis of Hankel radial modes and temporal frequencies $\Omega$, both operators are diagonal and therefore commute. 

%Concretely, \texttt{half\_linear} multiplies by the QDHT matrix, takes the FFT along time, multiplies by the single combined phase
%\begin{equation}
%\exp\!\left[i\left(\hat{D}(\Omega) - \frac{\rho^2}{2k_0}\hat{U}^{-1}(\Omega)\right)\frac{dz}{2}\right],
%\end{equation}
%applies the spectral mask, and then transforms back. Proceeding in this way costs one FFT and inverse FFT pair per call, instead of the two pairs that a separate frequency-resolved diffraction step would require on top of the dispersion step.

Inside the nonlinear part, the code distinguishes two phenomena that behave differently. The absorption coefficient $\alpha=\alpha_{ion}+\sigma_\omega N/2$, coming from photoionization loss and plasma absorption, only reduces the amplitude of $u$ without exchanging energy between the frequency components; it is therefore applied exactly as a multiplicative factor $u \to u\,e^{-\alpha\,dz/2}$ on each side of the rest of the nonlinear step. The remaining part, that is to say the Kerr effect, self-steepening, the dispersive part of the plasma response, and the STE polarizability, constitutes a genuine differential equation in $z$ and is integrated with a fourth-order Runge-Kutta scheme over the substep. Since the exponential factor already applies the zeroth order of the two dissipative terms, \texttt{split} adds $\alpha u$ back into the right-hand side it returns, so that the two treatments do not double-count the same loss.

The full sequence executed by \texttt{Integrator\dotb step} for one propagation step is therefore a half linear step, a half absorption step, the four Runge-Kutta stages, another half absorption step, and a final half linear step. A radial mask is applied at the end to damp anything reaching the edge of the box.

\subsection{Population Rate Equations on the GPU}

The rate equations are solved separately from the field envelope, by marching along the temporal axis at a fixed radius, once before every propagation step. The equations for $N$ and $N_{STE}$ are stiff, since the photoionization and avalanche rates can change by orders of magnitude within a single time sample, so a plain explicit Euler step would require an impractically small $\Delta t$. To avoid this, the code treats the coupled pair as a locally linear system over each time interval $[t_j,t_{j+1}]$,
\begin{equation}
\frac{dx}{dt} = S + Lx
\end{equation}
with the source $S$ and the rate $L$ evaluated from the trapezoidal average of the intensity-dependent rates at $t_j$ and $t_{j+1}$. This linear equation admits a closed-form solution,
\begin{equation}
x(t_j+\Delta t) = x(t_j)\,e^{L\Delta t} + S\,\Delta t\,\varphi_1(L\Delta t), \qquad \varphi_1(z) = \frac{e^z-1}{z}
\end{equation}
which is what the function \texttt{exact\_exp\_step} in the CUDA kernel implements, making use of a Taylor expansion for $\varphi_1$ when $L\Delta t$ is close to zero in order to avoid dividing zero by zero. This is an exact integrator for the frozen coefficient problem, therefore it does not accumulate truncation error from the time-stepping itself. The only error arises from the assumption that $S$ and $L$ are constant over each short interval, which represents a much better approximation than freezing the whole nonlinear right-hand side of a generic ordinary differential equation.

\subsection{The Main Propagation Loop}

Figure~\ref{fig:main_loop} gathers the three preceding subsections into the loop that \texttt{Integrator\dotb propagate} executes once per $z$ step.

\begin{figure}[h!]
    \centering
    \begin{tikzpicture}[
        node distance=7mm and 14mm,
        box/.style={draw, thick, rounded corners, align=center, minimum width=4.8cm, minimum height=9mm, fill=blue!4},
        deco/.style={draw, thick, rounded corners, align=center, minimum width=4.8cm, minimum height=9mm, fill=yellow!12},
        dec/.style={draw, thick, diamond, aspect=2.4, align=center, minimum width=3.4cm, inner sep=1pt},
        arr/.style={-stealth, thick}
    ]
        \node[box] (rate) {Update $N$, $N_{STE}$\\ \footnotesize CUDA rate kernel, one thread per radius};
        \node[box, below=of rate] (lin1) {Half linear step\\ \footnotesize QDHT and FFT, diffraction $+$ dispersion};
        \node[deco, below=of lin1] (abs1) {Absorption half-step\\ \footnotesize exact exponential};
        \node[box, below=of abs1] (rk4) {RK4 nonlinear substep\\ \footnotesize Kerr, self-steepening, plasma, STE};
        \node[deco, below=of rk4] (abs2) {Absorption half-step};
        \node[box, below=of abs2] (lin2) {Half linear step};
        \node[box, below=of lin2] (maskr) {Radial edge mask};
        \node[dec, below=of maskr] (test) {$z < z_{\text{end}}$?};
        \node[box, below=10mm of test] (save) {Save \texttt{result\dotb npz}};

        \draw[arr] (rate) -- (lin1);
        \draw[arr] (lin1) -- (abs1);
        \draw[arr] (abs1) -- (rk4);
        \draw[arr] (rk4) -- (abs2);
        \draw[arr] (abs2) -- (lin2);
        \draw[arr] (lin2) -- (maskr);
        \draw[arr] (maskr) -- (test);
        \draw[arr] (test) -- node[right] {no} (save);
        \draw[arr] (test.west) -- ++(-2.6,0) node[midway, above] {yes, $z \mathrel{+}= dz$} |- (rate.west);
    \end{tikzpicture}
    \caption{One iteration of the main propagation loop in \texttt{Integrator\dotb propagate}. Every $z$ step first advances the carrier populations along the full time axis at fixed $r$, then applies the symmetric Strang sequence to the field envelope. Diagnostics are recorded every \texttt{save\_stride} steps, omitted here for clarity, and the loop repeats until $z$ reaches the end of the propagation window.}
    \label{fig:main_loop}
\end{figure}

Let us note the order of the two blocks at the top. The populations are updated first, by using the field from the end of the previous step, and then the field is propagated by keeping these populations fixed. This constitutes in itself a splitting between the field and the carriers, on top of the Strang splitting inside the field update, and it remains accurate as long as $dz$ is sufficiently small so that neither changes much over one step.